\documentclass{article}

\usepackage{amsmath,amssymb}
\usepackage{xurl}
\usepackage[hidelinks]{hyperref}
\setlength{\emergencystretch}{2em}

\input{command}

\title{Resolution of the Choi-Type Positivity Scope (Wolf Example~3.1)}
\author{}
\date{2026-07-02, revised 2026-08-24}

\begin{document}
\maketitle
\gapnote{scope-restriction}{resolved}

\section{Source statement}

Let $\mathcal M_d=\MN{d}$, let
$D\colon\mathcal M_d\to\mathcal M_d$ be the projection onto the diagonal
matrices, and let $U_{k0}$ be the cyclic shift by $k$ coordinates.
Example~3.1 in Chapter~3 of Ref.~\cite{Wolf2012Quantum} introduces the
Choi-type map
\begin{align}
    T_C(X)
        &=(d-n)D(X)-X
          +\sum_{k=1}^{n}D(U_{k0}XU_{k0}^{\dagger}).
    \label{eq:wolf_ex3_1_choi_positivity_subcase_scope_choi_type_map}
\end{align}
This is Equation~(3.20) of Ref.~\cite{Wolf2012Quantum}.\footnote{Local source:
\path{Notes/WolfNoteTexSource/ch03_positive_not_completely.tex},
lines 357--365.}
The source states that $T_C$ is positive and indecomposable whenever
\begin{align}
    1
        &\leq n\leq d-2.
    \label{eq:wolf_ex3_1_choi_positivity_subcase_scope_source_parameter_range}
\end{align}
It imposes no hypothesis that the coordinates of an input vector be nonzero.

\section{Endpoint statements}

The endpoint arguments give independent proofs of two slices of the positivity
claim.

For the lower endpoint $n=1$, the development proves that $T_C$ maps every
positive semidefinite input to a positive semidefinite output in every
dimension $d\geq 3$. The scalar input is the boundary-inclusive cyclic
reciprocal inequality
\begin{align}
    \sum_{i\in\mathbb Z/d\mathbb Z}
        \frac{x_i}{(d-1)x_i+x_{i-1}}
        &\leq 1,
    \qquad x_i\geq 0,
    \label{eq:wolf_ex3_1_choi_positivity_subcase_scope_lower_endpoint_reciprocal}
\end{align}
where a summand with zero denominator is defined to be zero. In the interior
of the nonnegative cone, define
\begin{align}
    y_i
        &=\frac{x_{i-1}}{x_i},
    \label{eq:wolf_ex3_1_choi_positivity_subcase_scope_ratio_definition}\\
    \prod_{i\in\mathbb Z/d\mathbb Z}y_i
        &=1.
    \label{eq:wolf_ex3_1_choi_positivity_subcase_scope_ratio_product}
\end{align}
Clearing denominators in
(\ref{eq:wolf_ex3_1_choi_positivity_subcase_scope_lower_endpoint_reciprocal})
produces an expression in the elementary symmetric sums $e_k(y)$. The
arithmetic--geometric mean inequality gives
\begin{align}
    e_k(y)
        &\geq \binom{d}{k},
    \label{eq:wolf_ex3_1_choi_positivity_subcase_scope_symmetric_sum_bound}
\end{align}
and the resulting binomial expression vanishes for the constant family. This
is the proof route of Theorem~1 and Lemmas~2--3 of
Ref.~\cite[pp.~309--311]{TanahashiTomiyama1988Indecomposable}: Lemma~3 gives
the reciprocal estimate, and Theorem~1 applies the rank-one
diagonal-minus-projector criterion from Lemma~2.

The argument in
Ref.~\cite[pp.~309--311]{TanahashiTomiyama1988Indecomposable} assumes strictly
positive coordinates. If some $x_i$ vanishes, at most $d-1$ summands in
(\ref{eq:wolf_ex3_1_choi_positivity_subcase_scope_lower_endpoint_reciprocal})
are nonzero, and every nonzero summand is at most $1/(d-1)$. This proves the
boundary case as well.\footnote{Formal declarations:
\leanid{Matrix.choiType_cyclic_reciprocal_one_of_nonneg},
\leanid{Matrix.choiTypeMap_vecMulVec_posSemidef_one}, and
\leanid{Matrix.choiTypeMap_isPositiveMap_one} in
\doclink{QICLean/Channel/ChoiTypeMap/Positivity}.}
A spectral decomposition of a positive semidefinite matrix into rank-one
positive semidefinite summands then extends the rank-one result to positivity
of the map.

For the upper endpoint $n=d-2$, the development proves positivity on every
positive semidefinite input for all $d\geq 3$. This endpoint includes the case
$(d,n)=(3,1)$.\footnote{Formal declarations:
\leanid{Matrix.perm_reciprocal_sum_le_one},
\leanid{Matrix.choiTypeMap_vecMulVec_posSemidef_sub_two}, and
\leanid{Matrix.choiTypeMap_isPositiveMap_sub_two} in
\doclink{QICLean/Channel/ChoiTypeMap}.}
The backward window $x_{i-1},\ldots,x_{i-n}$ then omits exactly $x_i$ and
$x_{i+1}$. Define
\begin{align}
    T
        &=\sum_j x_j,
    \label{eq:wolf_ex3_1_choi_positivity_subcase_scope_total_mass}\\
    a_i
        &=T+x_i-x_{i+1}.
    \label{eq:wolf_ex3_1_choi_positivity_subcase_scope_upper_endpoint_weight}
\end{align}
The cyclic estimate becomes a special case of the permutation inequality
\begin{align}
    \sum_i\frac{x_i}{T+x_i-x_{\sigma(i)}}
        &\leq 1,
    \qquad x_i\geq 0,
    \label{eq:wolf_ex3_1_choi_positivity_subcase_scope_permutation_reciprocal}
\end{align}
where $\sigma$ is an arbitrary permutation.

Set $\delta_i=x_i-x_{\sigma(i)}$. The pair bound
$x_i+x_{\sigma(i)}\leq T$ gives
\begin{align}
    \frac{x_i}{T+\delta_i}
        &\leq
        \frac{x_iT-x_i\delta_i+\delta_i^2/2}{T^2}.
    \label{eq:wolf_ex3_1_choi_positivity_subcase_scope_permutation_summand_bound}
\end{align}
For a permutation,
\begin{align}
    \sum_i x_i\delta_i
        &=\frac{1}{2}\sum_i\delta_i^2.
    \label{eq:wolf_ex3_1_choi_positivity_subcase_scope_permutation_delta_identity}
\end{align}
Summing
(\ref{eq:wolf_ex3_1_choi_positivity_subcase_scope_permutation_summand_bound})
and using
(\ref{eq:wolf_ex3_1_choi_positivity_subcase_scope_permutation_delta_identity})
proves
(\ref{eq:wolf_ex3_1_choi_positivity_subcase_scope_permutation_reciprocal}).
This cyclic estimate is due to Ando, as recorded in the introduction of
Ref.~\cite{Yamagami1993Cyclic}.

\section{Resolution by Yamagami's variational argument}

The former scope restriction is resolved by the nonnegative cyclic reciprocal
estimate
\begin{align}
    \sum_{i\in\mathbb Z/d\mathbb Z}
        \frac{x_i}
        {(d-n)x_i+\sum_{k=1}^{n}x_{i-k}}
        &\leq 1,
    \qquad x_i\geq 0,\quad 2\leq n\leq d-3,
    \label{eq:wolf_ex3_1_choi_positivity_subcase_scope_middle_reciprocal}
\end{align}
again with zero-denominator summands defined to be zero. The corresponding
forward-window inequality is proved at the cardinal parameter. Index negation
converts it to
(\ref{eq:wolf_ex3_1_choi_positivity_subcase_scope_middle_reciprocal});
applying the result to $x_i=|v_i|^2$ proves positivity on rank-one positive
semidefinite inputs. Spectral decomposition then proves positivity of $T_C$
for every $d\geq 3$ and $1\leq n\leq d-2$.\footnote{Formal declarations:
\leanid{Yamagami.functional_card_le_one},
\leanid{Matrix.choiTypeRankOneWeight_reciprocal_sum_middle},
\leanid{Matrix.choiTypeMap_vecMulVec_posSemidef_middle}, and
\leanid{Matrix.choiTypeMap_isPositiveMap} in
\doclink{QICLean/Channel/ChoiTypeMap/YamagamiPositivity}.}

The middle range requires more than the endpoint arguments. The case $n=1$ is
Theorem~1 of
Ref.~\cite[pp.~309--311]{TanahashiTomiyama1988Indecomposable}, with
Lemmas~2--3 supplying the proof route used here. Yamagami proves the full
range $1\leq n\leq d-1$ by a variational argument
\cite{Yamagami1993Cyclic}. The argument combines Nowosad's uniqueness theorem
for critical points of the functional
\begin{align}
    x
        &\longmapsto
        \sum_i\frac{x_i}{(Sx)_i}
    \label{eq:wolf_ex3_1_choi_positivity_subcase_scope_variational_functional}
\end{align}
on the positive cone, a spectral condition for the Hessian at the constant
vector, and a projective-boundary analysis that counts vanishing summands
\cite{Nowosad1968Isoperimetric,Yamagami1993Cyclic}.

\subsection{The spectral and Hessian ingredients}

Let $C$ denote the cyclic shift matrix. For integers $d\geq 3$ and
$1\leq m\leq d-2$, and for $s\geq d$, define
\begin{align}
    S
        &=(s-m)I+C+\cdots+C^m.
    \label{eq:wolf_ex3_1_choi_positivity_subcase_scope_cyclic_matrix}
\end{align}
For every nontrivial standard character, the corresponding eigenvalue of $S$
lies in the open disk stated in Lemma~6 of
Ref.~\cite[pp.~523--524]{Yamagami1993Cyclic}.\footnote{Formal declaration:
\leanid{Yamagami.norm_symbol_sub_half_lt_half}.}
This is the corrected stride-one portion needed for Wolf's parameter range.

Lemma~6 of Ref.~\cite[pp.~523--524]{Yamagami1993Cyclic} is printed for
$1\leq m\leq d-l$ and $s\geq d$. At $l=1$, $m=d-1$, and $s=d$, however,
\begin{align}
    S
        &=I+C+\cdots+C^{d-1}
    \label{eq:wolf_ex3_1_choi_positivity_subcase_scope_false_endpoint_matrix}
\end{align}
has eigenvalue zero on every nontrivial character. The resulting disk bound is
an equality rather than the printed strict inequality. The endpoint becomes
valid when $s>d$, while $m\leq d-2$ is the exact range required here. The
false printed endpoint and its correction are recorded separately in
Ref.~\cite{gap:yamagami_lemma6_endpoint}.

Define the Hessian matrix
\begin{align}
    H
        &=s(S+S^{\mathsf T})-2S^{\mathsf T}S.
    \label{eq:wolf_ex3_1_choi_positivity_subcase_scope_hessian_matrix}
\end{align}
In the range $d\geq 3$, $1\leq m\leq d-2$, and $s\geq d$, the matrix $S$ has
nonnegative entries, every row and column sum is $s$, and $S$ is invertible.
Moreover,
\begin{align}
    H
        &\succeq 0,
    \label{eq:wolf_ex3_1_choi_positivity_subcase_scope_hessian_positive}\\
    \ker H
        &=\mathbb R\mathbf 1.
    \label{eq:wolf_ex3_1_choi_positivity_subcase_scope_hessian_kernel}
\end{align}
The proof diagonalizes $S$ and $H$ by the finite Fourier transform and proves
strict positivity away from the constant character.\footnote{The cyclic
matrix development is in
\doclink{QICLean/Analysis/YamagamiCyclicMatrix}. Formal declaration for the
off-character positivity statement:
\leanid{Yamagami.hessianSymbol_pos}.}

Let $f_{d,m,s}$ denote Yamagami's cyclic functional, and let
$\lambda_{S^{-1}}$ denote the corresponding inverse-operator functional. The
previously missing identity is
\begin{align}
    f_{d,m,s}(x)
        &=\lambda_{S^{-1}}(Sx).
    \label{eq:wolf_ex3_1_choi_positivity_subcase_scope_inverse_operator_identity}
\end{align}
It supplies an algebraic ingredient used in Lemmas~2--4 and Corollary~5 of
Ref.~\cite[pp.~522--523]{Yamagami1993Cyclic}, together with the corrected
range of Lemma~6 on pp.~523--524.\footnote{Formal declaration:
\leanid{Yamagami.functional_eq_lambdaT_inverseOperator_mulVec}.}

\subsection{The projective boundary}

The definitions retain Yamagami's forward window and use index negation to
obtain the backward Choi window. The simultaneous-singularity argument treats
a nonzero nonnegative boundary vector under the hypotheses
\begin{align}
    N
        &\geq 1,
    \label{eq:wolf_ex3_1_choi_positivity_subcase_scope_boundary_dimension}\\
    m
        &<N,
    \label{eq:wolf_ex3_1_choi_positivity_subcase_scope_boundary_window}\\
    N
        &\leq s.
    \label{eq:wolf_ex3_1_choi_positivity_subcase_scope_boundary_parameter}
\end{align}
A vanishing denominator produces a first positive coordinate after its zero
window. This makes $m$ summands tend to zero, and counting the complementary
summands gives
\begin{align}
    \frac{N-m}{s-m}
        &\leq \frac{N}{s}.
    \label{eq:wolf_ex3_1_choi_positivity_subcase_scope_boundary_count}
\end{align}
This is the global count indicated, but not printed, after Yamagami's
one-singularity calculation
\cite[pp.~524--525]{Yamagami1993Cyclic}. Its completion and attribution are
recorded in
Ref.~\cite{gap:yamagami_simultaneous_singularity_boundary}.

The regular-boundary recurrence in Ref.~\cite[p.~525]{Yamagami1993Cyclic} is
formalized for $N\geq 3$, $2\leq m\leq N-2$, and $s\geq N$, including reduced
vectors with singular denominators under Lean's totalized
convention.\footnote{Formal declaration:
\leanid{Yamagami.functional_le_deleteLastCoordinate}.}
Strong induction applies this recurrence at regular boundary points. If a
positive vector had functional value greater than one, normalization would
produce a compact superlevel closure in the standard simplex. The
simultaneous-singularity estimate excludes boundary points with a zero
denominator; continuity and the inductive estimate exclude the remaining
boundary points. Hence a maximizer must be strictly positive. The
Fourier--Hessian result and Nowosad's uniqueness theorem force that maximizer
to be scalar, where the functional equals one, giving a contradiction
\cite{Nowosad1968Isoperimetric,Yamagami1993Cyclic}.

\subsection{The finite-coordinate Nowosad argument}

The finite-coordinate input follows the source route in
Ref.~\cite[pp.~409--418]{Nowosad1968Isoperimetric}. Boundedness is proved for
the norm induced by the coordinate-sum functional, which is the finite
$L^2$ norm.\footnote{Formal declaration:
\leanid{Nowosad.exists_inducedHilbertNorm_bound}.}
Lagrange idempotents for the value classes of the generator show that the
coordinate derivation vanishes on the possibly proper Laurent algebra $P(w)$,
including the case in which coordinates of $w$ repeat.\footnote{Formal
declarations:
\leanid{Nowosad.mem_laurentSubalgebra_of_eq_on_valueClass} and
\leanid{Nowosad.coordinate_pointDerivation_eq_zero}.}
The genuine first and second variations yield the finite real specialization
of Nowosad's Theorem~1.8
\cite{Nowosad1968Isoperimetric}.\footnote{Formal declaration:
\leanid{Nowosad.finiteCoordinate_theorem_one_eight}.}
Its local-maximum corollary follows from
\begin{align}
    \lambda_{-T}
        &=-\lambda_T
    \label{eq:wolf_ex3_1_choi_positivity_subcase_scope_lambda_sign_identity}
\end{align}
rather than from reversing the orientation of Nowosad's theorem.

For Yamagami's matrix $S$, define the change of variables by
\begin{align}
    b
        &=\frac{Sa}{s},
    \label{eq:wolf_ex3_1_choi_positivity_subcase_scope_change_variable_b}\\
    x(t)
        &=sS^{-1}(b^t).
    \label{eq:wolf_ex3_1_choi_positivity_subcase_scope_change_variable_x}
\end{align}
The matrix $H$ remains the matrix defined in
(\ref{eq:wolf_ex3_1_choi_positivity_subcase_scope_hessian_matrix}). The
ordinary Euclidean Hessian satisfies
\begin{align}
    D^2f_S(\mathbf 1)[h,h]
        &=-s^{-3}\langle h,Hh\rangle.
    \label{eq:wolf_ex3_1_choi_positivity_subcase_scope_euclidean_hessian}
\end{align}
This identity fixes Yamagami's normalization. It is not formalized separately
as a Fr\'echet-derivative theorem; the development instead uses the
corresponding algebraic second-variation identity.\footnote{Formal declaration:
\leanid{Yamagami.minimumQuadraticForm_normalizedNegativeInverse_eq}.}

Under
(\ref{eq:wolf_ex3_1_choi_positivity_subcase_scope_hessian_positive}) and
(\ref{eq:wolf_ex3_1_choi_positivity_subcase_scope_hessian_kernel}), every
strictly positive local maximum is scalar.\footnote{Formal declaration:
\leanid{Yamagami.lemma_three_localMax_isScalar}.}
More precisely, a non-scalar maximum would produce the non-scalar tangent
\begin{align}
    sS^{-1}\log\left(\frac{Sa}{s}\right)
        &\in\ker H.
    \label{eq:wolf_ex3_1_choi_positivity_subcase_scope_nonscalar_tangent}
\end{align}
This is the uniqueness assertion in Lemma~3 of
Ref.~\cite[p.~522]{Yamagami1993Cyclic}; it does not assert that the unit vector
is itself a local maximum.\footnote{Formal declaration:
\leanid{Yamagami.pulledBackTangent_mem_hessianKernel_and_not_isScalarVector}.}

\subsection{Why the endpoint summand bound does not extend}

Two computations explain why elementary summand bounds do not cover the
remaining range. First, the cyclic estimate is tight in two distinct regimes.
At the constant vector, every summand equals $1/d$. For the geometric family
\begin{align}
    x_i
        &=t^i,
    \qquad i=0,\ldots,d-1,
    \label{eq:wolf_ex3_1_choi_positivity_subcase_scope_geometric_family}
\end{align}
as $t\to\infty$, each of the $d-n$ summands whose backward window avoids
wrap-around tends to $1/(d-n)$, while every remaining summand tends to zero.
The full sum again tends to one. Any summand-wise estimate would therefore
need to be sharp in both regimes.

Second, the proof for $n=d-2$ does not extend. Set
\begin{align}
    m
        &=d-n,
    \label{eq:wolf_ex3_1_choi_positivity_subcase_scope_complementary_window}\\
    \delta_i
        &=\sum_{k=1}^{m-1}(x_i-x_{i+k}),
    \label{eq:wolf_ex3_1_choi_positivity_subcase_scope_general_delta}\\
    a_i
        &=T+\delta_i.
    \label{eq:wolf_ex3_1_choi_positivity_subcase_scope_general_weight}
\end{align}
For $T>0$ and positive weights $a_i$, the reciprocal estimate is equivalent
to
\begin{align}
    \sum_i\frac{x_i\delta_i^2}{a_i}
        &\leq \sum_i x_i\delta_i.
    \label{eq:wolf_ex3_1_choi_positivity_subcase_scope_equivalent_delta_bound}
\end{align}
The endpoint proof factors this inequality through
\begin{align}
    \sum_i\frac{x_i\delta_i^2}{a_i}
        &\leq \frac{1}{m}\sum_i\delta_i^2,
    \label{eq:wolf_ex3_1_choi_positivity_subcase_scope_first_delta_step}\\
    \frac{1}{m}\sum_i\delta_i^2
        &\leq
        \frac{1}{2}\sum_{k=1}^{m-1}
            \sum_i(x_i-x_{i+k})^2,
    \label{eq:wolf_ex3_1_choi_positivity_subcase_scope_second_delta_step}\\
    \frac{1}{2}\sum_{k=1}^{m-1}
            \sum_i(x_i-x_{i+k})^2
        &=\sum_i x_i\delta_i.
    \label{eq:wolf_ex3_1_choi_positivity_subcase_scope_delta_energy_identity}
\end{align}
The first inequality follows from $a_i\geq mx_i$ and holds for every $m$. In
Fourier variables, the second inequality is
\begin{align}
    \left|\sum_{k=1}^{m-1}(1-\zeta^k)\right|^2
        &\leq
        \frac{m}{2}\sum_{k=1}^{m-1}|1-\zeta^k|^2
    \label{eq:wolf_ex3_1_choi_positivity_subcase_scope_fourier_endpoint_bound}
\end{align}
for every $d$th root of unity $\zeta$. It is an equality for $m=2$, but it
fails for $m=3$ and $\zeta=e^{i\pi/3}$:
\begin{align}
    \left|\sum_{k=1}^{2}(1-\zeta^k)\right|^2
        &=7,
    \label{eq:wolf_ex3_1_choi_positivity_subcase_scope_fourier_counterexample_left}\\
    \frac{3}{2}\sum_{k=1}^{2}|1-\zeta^k|^2
        &=6.
    \label{eq:wolf_ex3_1_choi_positivity_subcase_scope_fourier_counterexample_right}
\end{align}
This obstruction agrees with the historical development: the case
$(d,n)=(5,2)$ required a separate argument
\cite{Osaka1991Indecomposable}, and Yamagami later established the full range
by the variational proof \cite{Yamagami1993Cyclic}.

The formalization therefore follows the variational route and resolves
\href{https://github.com/LionSR/QICLean/issues/19}{QICLean issue \#19}. The
indecomposability assertion in the same paragraph of
Ref.~\cite{Wolf2012Quantum} is separate and remains tracked by
\href{https://github.com/LionSR/QICLean/issues/18}{QICLean issue \#18}.

\section{Classification}

The scope restriction is resolved. The formal positivity theorem has exactly
Wolf's hypotheses
\begin{align}
    d
        &\geq 3,
    \label{eq:wolf_ex3_1_choi_positivity_subcase_scope_resolved_dimension}\\
    1
        &\leq n\leq d-2,
    \label{eq:wolf_ex3_1_choi_positivity_subcase_scope_resolved_parameter_range}
\end{align}
and concludes that the Choi-type map
(\ref{eq:wolf_ex3_1_choi_positivity_subcase_scope_choi_type_map}) is
positive.\footnote{Formal declaration:
\leanid{Matrix.choiTypeMap_isPositiveMap}.}
The former endpoint-only restriction has been removed. Wolf's independent
indecomposability assertion remains tracked by QICLean issue~\#18.

\bibliographystyle{alpha}
\bibliography{references}

\end{document}
